\documentclass[11pt]{article}

\usepackage[margin=1in]{geometry}
\usepackage{amsmath,amssymb}
\usepackage{hyperref}
\usepackage{enumitem}

\hypersetup{
  colorlinks=true,
  linkcolor=blue,
  urlcolor=blue,
  citecolor=blue
}
\emergencystretch=2em

\newcommand{\Free}{\mathcal{F}}
\newcommand{\R}{\mathbb{R}}

\title{A Literature-Implied Answer to the Schur/RNP Question for Lipschitz-Free Spaces}
\author{Math discovery packet}
\date{June 23, 2026}

\begin{document}
\maketitle

\section{Status}

This is a lightweight \emph{literature-implied answer}, not an original
solution.  The source paper is M. Cuth, M. Doucha and P. Wojtaszczyk,
\emph{Embeddings of Lipschitz-free spaces into \(\ell_1\)},
arXiv:1909.05285.  The supporting paper is R. J. Aliaga, C. Gartland,
C. Petitjean and A. Prochazka, \emph{Purely 1-unrectifiable metric spaces
and locally flat Lipschitz functions}, arXiv:2103.09370.

\section{Source Question}

The source paper first asks, in Question~1 of the introduction, whether every
closed negligible subset \(M\) of an \(\R\)-tree has \(\Free(M)\)
isometrically embeddable into some \(\ell_1(\Gamma)\).  That question is
answered negatively inside the same paper by Theorem~1, so it is not recorded
as a separate literature answer.

The surviving external question appears after Corollary~2.4
(\texttt{c:RNPSchurL1}).  There the source paper proves that, for a closed
subset \(M\) of an \(\R\)-tree, the following are equivalent:
\begin{enumerate}[leftmargin=*]
\item \(M\) is negligible,
\item \(\Free(M)\) has the Schur property,
\item \(\Free(M)\) has the Radon--Nikodym property,
\item \(\Free(M)\) contains no isomorphic copy of \(L_1\).
\end{enumerate}
It then states that it is unknown whether any of the corresponding
equivalences, in particular Schur \(\Longleftrightarrow\) Radon--Nikodym, hold
for general Lipschitz-free spaces.

\section{Supporting Theorem}

ArXiv:2103.09370 proves the broader theorem.  Theorem~C in its introduction,
and Theorem~4.6 (\texttt{thm:equivalences}) in the body, state that for every
metric space \(M\) the following conditions are equivalent:
\begin{enumerate}[leftmargin=*]
\item the completion of \(M\) is purely \(1\)-unrectifiable,
\item \(\Free(M)\) has the Radon--Nikodym property,
\item \(\Free(M)\) has the Krein--Milman property,
\item \(\Free(M)\) has the Schur property,
\item \(\Free(M)\) contains no isomorphic copy of \(L_1\).
\end{enumerate}
The body theorem also includes the equivalent quantitative formulation that
\(\Free(M)\) contains no \((1+\varepsilon)\)-isomorphic copy of \(L_1\) for
some \(\varepsilon>0\).

\section{Identification}

The free space over a metric space and over its completion are canonically the
same for the purposes of the theorem, as noted at the start of the proof of
Theorem~4.6 in the supporting paper.  Therefore the implications
\[
  \Free(M)\text{ has the Schur property}
  \quad\Longleftrightarrow\quad
  \Free(M)\text{ has the Radon--Nikodym property}
\]
hold for every Lipschitz-free space.  This gives an affirmative answer to the
general Schur/RNP equivalence question left open after Corollary~2.4 of
arXiv:1909.05285.

\section{Provenance and Scope}

This packet is placed in \texttt{literature\_implied\_answers}.  The
supporting paper explicitly treats the broad Schur/RNP problem for
Lipschitz-free spaces and says the equivalence was suggested by previous
work, but a source search did not find an explicit citation to arXiv:1909.05285
or to the exact sentence in that paper.  Hence the relation is recorded as an
agent-identified consequence of a decisive later theorem.

No part of this packet claims a new proof.  Its role is duplicate memory:
future agents should not attack the Schur/RNP equivalence for all
Lipschitz-free spaces as open.

\section{References}

\begin{itemize}[leftmargin=*]
\item M. Cuth, M. Doucha and P. Wojtaszczyk, \emph{Embeddings of
Lipschitz-free spaces into \(\ell_1\)}, arXiv:1909.05285.
\item R. J. Aliaga, C. Gartland, C. Petitjean and A. Prochazka, \emph{Purely
1-unrectifiable metric spaces and locally flat Lipschitz functions},
arXiv:2103.09370.
\end{itemize}

\end{document}
